% page 262 du fac-simile (image p-261 du scan)
% bloc 160.0 x 237.7 mm ; marge gauche 8.0 mm
\pgc{-12.700mm}{0.000mm}{
\l{\cel{3}254}
\l{}
\l{\cel{1}\sou{kalefaktar}. (\sou{netrans.})\cel{2}(\sou{fiziko)}.\cel{2}Aquo-guteti, disjetita}
\l{\cel{3}sur metala plako kelke varmigita; vaporeskas tre rapide;}
\l{\cel{3}ma se la surfaco esas tre varma, nula ebulio eventas, e}
\l{\cel{3}la guteti \textquotedbl{}kalefaktas\textquotedbl{}, t.e. li inter-asemblas e formacas}
\l{\cel{3}globeto sfera. - DEFIS.}
\l{}
\l{\sou{kalemburo}. Vorto-ludo fondata sur la senco-difero inter vorti}
\l{\cel{3}qui audale aspektas inter-sama. - DFR.}
\l{}
\l{\sou{kalendo}. Nomo donita da la Romani antiqua a la dio unesma di}
\l{\cel{3}singla monato. - DEFIRS.}
\l{}
\l{\sou{kalendario}.\cel{3}I. Sistemo di divido di tempo, per yari, monati,}
\l{\cel{3}e dii. - II.Tabelo ek la dii, monati, sezoni, pri singla}
\l{\cel{3}yaro, qua generale indikas, pri singla dio, la nomo disanto.}
\l{\cel{3}- DEFIRS.}
\l{}
\l{\sou{kalendulo}.\cel{2}(\sou{bot.}) Planto ek la familio \textquotedbl{}kompozaji\textquotedbl{}, di qua}
\l{\cel{3}la kapituli esas flava e radio-dispozita.- L.\cel{1}\sou{calendula}.}
\l{}
\l{\sou{kalenturo}. (\sou{patol.)}\cel{2}Kazo di febro ardorigiva, kun deliro,}
\l{\cel{3}qua eventas ulafoye che la navani kande la navo trairas}
\l{\cel{3}la zono torida. -DEFIS.}
\l{}
\l{\sou{kalesho.}\cel{2}Veturo eleganta,quar-rota, maxim-multa-kaze apertita}
\l{\cel{3}avane, e provizita (dope) ye tendo ek ledro qua esas des-}
\l{\cel{3}faldebla e retrofaldebla segun-vole. - DEFIR.}
\l{}
\l{\sou{kalfatar}. (\sou{trans.}) Igar aquo-nepermeebla (barko, ramo-batelo),}
\l{\cel{3}per stopar la junturi, la fenduri, la trui, per stupo gudro-}
\l{\cel{3}imprenita. - DFIS.}
\l{}
\l{\sou{kalio}.\cel{2}(\sou{kemio)}. Korpo simpla, metala, volatila, qua oxideskas}
\l{\cel{3}per kontaktar kun aero humida. -\cel{1}\sou{Simbolo kemiala}\cel{1}: K.\cel{2}- D.}
\l{}
\l{\sou{kalibro.}\cel{2}Modelo sur qua esas trasita la konturi, la dimen-}
\l{\cel{3}sioni, di la objekti fabrikenda. - DEFIRS.}
\l{}
\l{\sou{kalic}o.I. Vazo sakra en qua la sacerdoto sakrigas la vino di}
\l{\cel{3}la eukaristio.- II. Envelopilo qua kovras la infra parto}
\l{\cel{3}di la flor-korolo. - EFIS.}
\l{}
\l{\sou{kalidoskopo}.\cel{2}Aparato qua kontenas mikra objekti multa-kolora}
\l{\cel{3}qui susituas diverse ed intersimetrie kande onu agitas la}
\l{\cel{3}instrumento. - DEFIRS.}
\l{}
\l{\sou{kalifo}. Suvereno Mohamedista, qua unionas en su la povo po-}
\l{\cel{3}litikala e la povo religiala. - DEFIRS.}
\l{}
\l{\sou{kaligrafar}. (trans.)\cel{2}Skribar (ulo) per literi eleganta, or-}
\l{\cel{3}nita. - DEFS.}
\l{}
\l{\sou{kaliko}.\cel{2}Telo de kotono, min delikata kam perkalo. - DEFIRS.}
}
